\chapter{A nossa responsabilidade}\label{ch:g9:our-responsibility}

Nove anos deste livro terminam aqui, e os fios dele se
encontram num fato simples: a \omterm{def:g5:biodiversity:species}{espécie} capaz de ler esta
página é a única \omterm{def:g5:biodiversity:species}{espécie} que entende a maquinaria a que
pertence --- e entender, como o \cref{ch:g8:contraception}
disse primeiro, traz escolha. Este capítulo de encerramento
reúne as responsabilidades que o entendimento cria: com a
própria saúde, com a saúde das outras pessoas e com a teia de
\omterm{def:g6:exploring-our-environment:components}{seres vivos} que carrega todos nós.

\section{A responsabilidade com a própria maquinaria}

\begin{proposition}[A saúde como autogestão informada]\label{prop:g9:our-responsibility:self}
Os fios de saúde do livro, amarrados:
\begin{itemize}
  \item o contrato diário
        (\cref{prop:g5:healthy-choices:contract}) --- comida,
        movimento, \omterm{prop:g1:healthy-habits:needs}{sono}, limpeza --- é manutenção de órgãos
        que você já sabe justificar sistema por sistema
        (\cref{prop:g7:heart-and-circulation:keep},
        \cref{met:g7:lungs-and-gas-exchange:keep},
        \cref{met:g8:protecting-the-nervous-system:manual});
  \item os mecanismos das substâncias são conhecidos
        (\cref{prop:g8:protecting-the-nervous-system:substances}),
        e a assimetria da armadilha tem preço
        (\cref{rem:g5:healthy-choices:never});
  \item as escolhas reprodutivas têm o mapa e os endereços
        delas (\cref{ch:g8:contraception});
  \item e a história do corpo é em parte distribuída
        (\cref{ch:g9:heredity-and-traits}): conhecer as
        tendências de uma família --- os cabelos brancos
        precoces, os \omterm{def:g4:heart-and-blood:heart}{corações} que entopem --- não é destino,
        e sim um cronograma de manutenção
        (\cref{met:g9:unique-individuals:claims}: faixas,
        preenchidas pela vida).
\end{itemize}
\end{proposition}
\admitted

\begin{example}[A consulta de rotina, lida como biologia]\label{ex:g9:our-responsibility:checkup}
Uma consulta de rotina é este livro dentro de um atendimento:
o \omterm{prop:g4:heart-and-blood:pulse}{pulso} e a pressão (a rede lida à beira da estrada,
\cref{ex:g7:heart-and-circulation:measures}), o martelinho de
reflexo (\cref{pb:g8:nervous-communication:1}), a caderneta
de \omterm{prop:g9:immune-defenses:vaccination}{vacinação} (\cref{pb:g9:immune-defenses:1}), a história de
família (as tendências distribuídas) e os conselhos de sempre
--- o contrato, de novo. A medicina é em boa parte manutenção
aplicada de mecanismos cujos manuais agora são seus.
\end{example}

\section{A responsabilidade com os outros}

\begin{proposition}[A saúde é maquinaria compartilhada]\label{prop:g9:our-responsibility:others}
Três mecanismos fazem da saúde uma conta compartilhada:
\begin{itemize}
  \item a \emph{transmissão}
        (\cref{def:g9:microbes-and-infection:stages}): a sua
        \omterm{def:g1:healthy-habits:hygiene}{higiene}, o modo como você cobre a tosse e o seu dia
        de doente em casa são proteção para as outras
        pessoas;
  \item a \emph{\omterm{def:g9:immune-defenses:memory}{imunidade} reunida}
        (\cref{prop:g9:immune-defenses:vaccination}): a
        \omterm{prop:g9:immune-defenses:vaccination}{vacinação} protege, através de você, quem não pode ser
        treinado;
  \item os \emph{bens comuns da medicina}: os antibióticos
        mantidos funcionando
        (\cref{exo:g9:how-species-change:10}) e o \omterm{prop:g4:journey-of-food:absorb}{sangue}
        doado (\cref{ex:g9:unique-individuals:medicine} --- o
        registro funciona com doadores) são recursos
        compartilhados, gastos ou estocados por escolhas
        individuais.
\end{itemize}
\end{proposition}
\admitted

\begin{example}[A conta compartilhada de um adolescente]\label{ex:g9:our-responsibility:account}
Os lançamentos comuns da semana: ficou em casa com febre (a
transmissão cortada); reforços de vacina em dia (o fundo
comum sustentado); um tratamento com antibiótico terminado na
primavera (os bens comuns não gastos); uma primeira doação de
\omterm{prop:g4:journey-of-food:absorb}{sangue} planejada para os dezoito anos (o registro estocado).
Nada heroico --- a lição do \cref{ex:g5:healthy-choices:day}
na escala de uma cidade: a saúde pública são dias comuns, bem
cumpridos.
\end{example}

\section{A responsabilidade com o mundo vivo}

\begin{proposition}[O cuidado, reenunciado com todas as razões dele]\label{prop:g9:our-responsibility:stewardship}
Os fios da natureza, amarrados com os mecanismos deles:
\begin{itemize}
  \item toda \omterm{def:g3:food-chains:chain}{cadeia alimentar} começa no verde
        (\cref{prop:g6:plants-are-producers:firstlink}), e
        toda mesa é a produção de um ecossistema
        (\cref{prop:g6:producing-our-food:costs});
  \item as teias absorvem choques pela variedade
        (\cref{prop:g5:biodiversity:web}), e a variedade é a
        reserva com que a seleção funciona
        (\cref{prop:g9:unique-individuals:reserve},
        \cref{exo:g9:how-species-change:12}) --- a
        \omterm{def:g5:biodiversity:biodiversity}{biodiversidade} não é paisagem, e sim capital de giro;
  \item a extinção fecha a porta para sempre
        (\cref{prop:g5:biodiversity:loss}) --- a única
        exceção deliberada, a varíola, prova como a palavra é
        absoluta (\cref{exo:g9:immune-defenses:12});
  \item as regras do uso duradouro são conhecidas
        (\cref{prop:g5:people-and-nature:lasting}), as
        violações delas são diagnosticáveis
        (\cref{ex:g7:oxygen-in-the-water:waste}) e os reparos
        delas estão demonstrados
        (\cref{ex:g5:people-and-nature:river},
        \cref{ex:g5:biodiversity:storks}).
\end{itemize}
O cuidado é a mesma mentalidade de manutenção da saúde,
aplicada ao corpo maior dentro do qual vivemos.
\end{proposition}
\admitted

\begin{example}[Para que serve a biologia de um cidadão]\label{ex:g9:our-responsibility:citizen}
Os votos, as compras e os hábitos da vida adulta chegam sem
parar vestidos de perguntas para as quais este livro equipou:
esta pescaria está tirando mais depressa do que a renovação?
Esta política protege o habitat dos polinizadores ou o
asfalta? Esta dieta milagrosa, esta cura ou este ``\omterm{def:g9:genes-and-dna:gene}{gene} do
X'' que anunciam são coerentes
(\cref{met:g9:unique-individuals:claims},
\cref{met:g8:contraception:evaluate})? Este antibiótico
deveria estar neste cocho? Um cidadão capaz de fazer essas
auditorias é para o que serviram nove anos de capítulos.
\end{example}

\begin{method}[O método de quem se forma]\label{met:g9:our-responsibility:method}
Ao sair do Livro 1, leve o método mais do que os fatos:
\begin{enumerate}
  \item observe antes de concluir (o
        \cref{met:g1:animals-around-us:observe} começou; e
        todo censo e todo caderno desde então);
  \item teste com justiça --- uma diferença só, e um controle
        (\cref{met:g4:what-plants-need:fairtest}, usado em
        todos os anos seguintes);
  \item nomeie o mecanismo --- nenhuma alegação está auditada
        enquanto o elo dela não for nomeado
        (\cref{met:g8:contraception:evaluate});
  \item pese as testemunhas independentes
        (\cref{met:g9:evidence-of-evolution:weigh});
  \item e respeite a maquinaria viva, dentro do seu corpo e
        fora dele: ela é antiga, sutil, reparável dentro de
        certos limites --- e cabe a você conservá-la.
\end{enumerate}
\end{method}

\begin{remark}[O que vem a seguir]\label{rem:g9:our-responsibility:next}
O volume do ensino médio retoma esses fios no aumento
seguinte: a maquinaria molecular da \omterm{def:g5:first-look-at-cells:cell}{célula}, a genética com a
matemática completa dela, os sistemas imunológico e nervoso a
sério, a ecologia como ciência quantitativa e a \omterm{prop:g9:how-species-change:selection}{evolução} como
o quadro geral que ela é. As perguntas serão as mesmas,
feitas mais de perto. Traga o método.
\end{remark}

\section{Exercícios}

\begin{exercise}[$\star$]\label{exo:g9:our-responsibility:1}
Liste os quatro fios da responsabilidade consigo mesmo na
\cref{prop:g9:our-responsibility:self}.
\end{exercise}

\begin{exercise}[$\star$]\label{exo:g9:our-responsibility:2}
Nomeie os três mecanismos que fazem da saúde algo
compartilhado, cada um com o capítulo dele.
\end{exercise}

\begin{exercise}[$\star$]\label{exo:g9:our-responsibility:3}
Por que uma tendência de família conhecida é um cronograma, e
não um destino?
\end{exercise}

\begin{exercise}[$\star$]\label{exo:g9:our-responsibility:4}
Leia três instrumentos da consulta de rotina como os
capítulos deles.
\end{exercise}

\begin{exercise}[$\star$]\label{exo:g9:our-responsibility:5}
Por que a \omterm{def:g5:biodiversity:biodiversity}{biodiversidade} é ``capital de giro'', e não
paisagem? Dois mecanismos.
\end{exercise}

\begin{exercise}[$\star$]\label{exo:g9:our-responsibility:6}
Enuncie de memória os cinco passos do método de quem se
forma.
\end{exercise}

\begin{exercise}[$\star\star$]\label{exo:g9:our-responsibility:7}
Justifique o dia de doente em casa como um ato de saúde
pública, com a etapa nomeada.
\end{exercise}

\begin{exercise}[$\star\star$]\label{exo:g9:our-responsibility:8}
De quem é o escudo que a sua \omterm{prop:g9:immune-defenses:vaccination}{vacinação} faz? Responda com o
fundo comum e nomeie dois dependentes.
\end{exercise}

\begin{exercise}[$\star\star$]\label{exo:g9:our-responsibility:9}
Audite como faria o \cref{ex:g9:our-responsibility:citizen}:
``o nosso spray novo mata todo inseto do pomar --- proteção
total''.
\end{exercise}

\begin{exercise}[$\star\star$]\label{exo:g9:our-responsibility:10}
Explique ``os antibióticos são um bem comum'' com a seleção da
enfermaria --- e os dois hábitos que mantêm esse bem comum
estocado.
\end{exercise}

\begin{exercise}[$\star\star$]\label{exo:g9:our-responsibility:11}
Aplique os cinco passos do método a uma publicação viral de
``cura milagrosa'', um passo de cada vez.
\end{exercise}

\begin{exercise}[$\star\star\star$]\label{exo:g9:our-responsibility:12}
Escreva você mesmo o parágrafo de encerramento do livro: uma
frase para a \omterm{def:g5:first-look-at-cells:cell}{célula}, uma para o corpo, uma para a teia, uma
para a árvore genealógica --- e o lugar do leitor nas quatro.
\end{exercise}

\section{Problema: A auditoria de formatura}

\begin{problem}[{Problema de fim de semana --- um ano de uma
cidade, auditado com o livro inteiro}]\label{pb:g9:our-responsibility:1}
O ano (inventado) da sua cidade em manchetes: uma onda de
gripe no inverno; um debate na primavera sobre pulverizar os
plátanos; uma mortandade de \omterm{prop:g3:sorting-living-things:groups}{peixes} no lago do parque no
verão; uma campanha no outono pela nova reserva do brejo; e,
o ano inteiro, as campanhas de \omterm{prop:g9:immune-defenses:vaccination}{vacinação} e de doação do
posto.

\medskip\noindent\textbf{Parte I --- O arquivo da saúde.}
\begin{enumerate}
  \item A onda de gripe: nomeie as vias de transmissão que os
        conselhos da cidade miraram, e por que o posto
        vacinou primeiro a equipe da casa de repouso.
  \item As campanhas: diga o que o registro de \omterm{prop:g4:journey-of-food:absorb}{sangue} e a
        caderneta de \omterm{prop:g9:immune-defenses:vaccination}{vacinação} estocam cada um, e quem saca
        deles.
  \item Um colunista zomba do ``escarcéu com lavar as mãos''.
        Responda com a enfermaria do
        \cref{ex:g9:microbes-and-infection:history} --- um
        parágrafo.
  \item A farmácia relata a volta de tratamentos com
        antibiótico não terminados. Que bem comum está sendo
        gasto, e por que mecanismo?
\end{enumerate}

\medskip\noindent\textbf{Parte II --- O arquivo da natureza.}
\begin{enumerate}[resume]
  \item O debate da pulverização: faça a auditoria do pomar
        (\cref{exo:g9:our-responsibility:9}) sobre o ``matar
        todo inseto'' --- quem mais morre, e que serviços
        gratuitos param?
  \item A mortandade de \omterm{prop:g3:sorting-living-things:groups}{peixes}: uma semana quente e parada e
        um lago verde-ervilha. Diagnostique com o
        \cref{ch:g7:oxygen-in-the-water} e receite.
  \item A campanha do brejo: defenda a reserva com três
        mecanismos --- o habitat, a teia e a reserva de
        variedade --- uma frase para cada um.
  \item Um proprietário de terras pergunta ``o que o brejo já
        fez pela cidade?''. Responda com o livro-caixa do
        trabalho gratuito (\cref{exo:g5:biodiversity:11}, o
        armazenamento de cheias do rio do
        \cref{ex:g5:people-and-nature:river}).
\end{enumerate}

\medskip\noindent\textbf{Parte III --- O encerramento de quem
se forma.}
\begin{enumerate}[resume]
  \item Escolha a melhor e a pior decisão do ano pelas regras
        do uso duradouro, e justifique cada veredito.
  \item Mostre os cinco passos do método em ação no jeito
        como a mortandade de \omterm{prop:g3:sorting-living-things:groups}{peixes} foi resolvida --- da
        primeira observação ao mecanismo aceito.
  \item Redija o editorial de fim de ano do jornal da cidade
        em quatro frases: a saúde compartilhada, a natureza
        guardada, as alegações auditadas, o método mantido.
  \item A última linha da última página do Livro 1 é sua para
        escrever: uma frase, dirigida ao leitor que começou
        no caracol do \cref{ch:g1:living-or-not}.
\end{enumerate}
\end{problem}
